\chapter[Sermon 53. The Description of the Conflict of Christ and the Church with the Dragon: the Dragon Is Overcome, the Heavenly Dwellers Sing Praises.]{The Description of the Conflict of Christ and the Church with the Dragon: the Dragon Is Overcome, the Heavenly Dwellers Sing Praises.}
\label{sermon:053}
\markright{Sermon 53. The Description of the Conflict of Christ and the ...}

And there was a great battle in heaven; Michael and his angels fought with the dragon, and the dragon fought with his angels, but they prevailed not, and was cast out of heaven. And the dragon, that old serpent, who is called the devil and Satan, was cast out, he who deceives the whole world. And he was cast to the earth, and his angels were cast out with him also. And I heard a loud voice in heaven, saying, “Now has salvation and power and the kingdom come to our God, and the power of the Lord is Christ’s, because he has been cast down, he who accuses them before God day and night.” And they overcame him by the blood of the Lamb and by the word of their testimony, and they did not love their lives to the point of death. Therefore, rejoice, heavens and those who dwell therein. Woe to the inhabitants of the earth and the sea, for the devil has come down to you, full of great wrath, because he knows that he has but a short time.

The Apostle has spoken of the parts of the remarkable conflict and worthy battle; he has spoken also of attempts and purposes of the dragon, which indeed directs all his schemes to this intent, that he may devour all godliness, that is, destroy it utterly. He has shown how he began to move war against the church, which fled into the wilderness; and now, as it were leaving the woman in the wilderness, he seems to bring forth other soldiers, who give battle to the dragon, and most valiantly do challenge and also defeat him and all his power. John therefore describes the singular fight of one most excellent, namely Michael, who overcame the dragon; and describes the general fight joined with that particular. For he adds that all the angels of Michael fought against the dragon.

First, it is shown to be the place of the fight or conflict. For he says, “A great battle was fought in heaven.” And it is evident that Satan was, at the beginning of all things, cast out of heaven into the earth, and therefore he does not wage war in heaven, nor raise any tumult there. For heaven is a place of rest and joy, not of debate and contention. Therefore, this must be attributed to the vision. For the Lord has, in heaven, represented this battle by signs to be seen, which in reality is fought on earth in the midst of the church.

Here is set forth an image of a significant battle, showing what has been, and what is yet to be done on earth. I said even now that this vision was indeed specific, but to have a general battle appended. For Michael fights, as captain of this war, and Michael’s angels fight also. This must be carefully discerned, although Michael and his angels make up only one part. On the other side fights the dragon, as emperor of this war, and his angels fight also. And these truly make no other parties than we have heard before at the beginning of this chapter: that the parties in this fight are the church and the devil. Nevertheless, lest the victory should be attributed to the church, and not rather to Christ, the woman must now be omitted, and Michael brought in fighting. There is some difficulty in these things, but it will be easy enough for him who will mark everything in order.

First we must see what Michael is, and there is no doubt that the angel Michael appeared in the vision, with an army of angels fighting. And on the contrary part, the Dragon fought with a host of devils. But since we heard at the beginning that these were signs, they must needs signify and betoken other things. I suppose, therefore, that Christ is signified here, the head of his church, king and protector, with his members, Apostles, Martyrs, and faithful. It is not rare that Christ should be figured to us by angels; indeed, it is most customary that angels are called the ambassadors of God and the faithful servants of Jesus Christ. Therefore, Christ, the head of the church, and the faithful members of Christ, fight against the Dragon, yet in a different way. For Christ overcame him alone in the combat without the help of any creature, while in temptations he discomfited him at last, and also by dying on the cross and rising again from the dead, he utterly broke his head. This is the only, true, and singular victory, whereby afterwards are obtained the victories of Christ's members, derived from that general fight, wherein Christ fights not now only hand to hand with the Devil, but all the members of Christ at all times, under Christ their Captain, fight against the Devil, and in the virtue or victory of Christ, fight and overcome, as we shall hear by and by in the song of praise.

\index{Daniel (Book of)}\index{Ezekiel (Book of)}But for great and sundry causes we affirm that Christ is figured and signified to us under the type of Michael. We know, by the scriptures as many of us as are learned, that Michael, as also Gabriel, are the names of good angels of God. Michael signifies, “Who is as God?” And who, I pray you, is such as God, but in whom the express image of the Father’s substance, which is the invisible Image and word of the Father from the beginning, I mean the very Son of God Jesus Christ? Michael in the tenth and twelfth chapters of Daniel, is president, protector, and patron of the Jewish nation. And it is plain that the people of Israel had from the beginning no other tutor and patron but Messiah himself, the blessed Seed. This appears in the seventh chapter of Isaiah, where we read that the Lord spared the people of Judah and the princely city for Christ. In another place he says most openly, “I will defend that city for myself and for my servant David.” And David is called Christ in the thirty-fourth chapter of Ezekiel. Christ is therefore in very deed governor of his people, which nevertheless in defending and delivering his people, uses the ministry of angels; who also attribute nothing to themselves, but all glory to God alone. Moreover, that excellent victory cannot without offense of godliness be ascribed to Michael the archangel. For so, omitting our Messiah Christ, we should commend angels being made and worthy to be called angelical, rather than Christians. In the law was written, “The seed of the woman shall bruise the serpent’s head.” But the Lord never took the nature of an angel, but the seed of Abraham, and by sin condemned sin. There shall follow anon in the song, “Now is salvation and power, \&c.” And it is added, “for the Devil is cast out.” And this salvation has Christ alone accomplished, wherefore it is necessary that Christ the conqueror of Satan be signified by Michael.

And the dragon fought hand to hand against the Lord, not only contending with him in the wilderness, but also never ceasing to tempt and assail him as long as he lived on earth. He stirred up against him the Pharisees and princes of the people, kings and the Roman governor, and thus at last broke the Lord’s heel. This was the greatest fight of the dragon. The same dragon now inspires kings and princes, wicked priests and cruel men, his agents who may wage war upon the church. And all these truly do persecute and vex the church in the power of the red dragon. Histories declare the same to have been done before Christ’s time; the same testify, and experience proves, that the like has been done from the ascension of Christ into heaven, unto this present day, and unto the world’s end.

\index{Antichrist}Now it is also declared with what fortune they fought on either side: namely, most fortunately concerning Christ, and most unfortunately as touching the Devil or red Dragon. And in this fight, as also in the song immediately following, is contained the whole fruit of this disputation. For from this, all the faithful may learn that Satan, our enemy, is unarmed; and that Christ in this conflict is on our side, as our Emperor and captain at all times, by whom all the faithful may easily overcome in all conflicts. Therefore, this matter of battle and victory is placed immediately after the beginning of the most dangerous battle with Antichrist and Antichristians, who are the brood or offspring and scales of the serpent, and champions of the Dragon, for comfort and consolation. And the natural order is here altered, which does not treat of the success of the battle until he has set forth all the conflict before. But this battle shall be continued hereafter in the rest of the twelfth and all of the thirteenth chapter.

He states this in three phrases: first, the victory of Christ; secondly, the victory of all Christians. The first is, “they prevailed not, they had no strength.” Doubtless the devil’s power is great if God permits, and clearly greatest, in consideration of God’s just judgment, as also appears in Job, that he is able to quench and break the strongest things. But the Lord says in the gospel: “The prince of this world came, and against me he has nothing.” Again in the gospel: “The gates of Hell shall not prevail against her, the rock I mean, and secondly against the church.” Although therefore the devil makes a horrible uproar and cruelly rages against Christ and his church, he is without force. For the virtue of Christ prevails.

\index{John, Saint}The second part is, neither is their place any longer found in heaven, which manner of speech signifies nothing other than that the condemned angel is removed from all dignity, glory, and power; moreover, that he has no more place in the church, or among the elect of God; not that the devil should not return, or should not tempt, or renew war, but because he has no permanent place. This relates to what the Lord so often repeated in the Gospel, and now the prince of this world is cast out, in the 12th and 16th chapters of Saint John. Moreover, by other passages of Scripture it is manifest that the devil is shut out of heaven. And it shall be easy for us to shut him out, who, being cast out by the Son of God, has no place in us, unless we ourselves give place to him. Which we should not do, the Lord admonishes us diligently, that we should watch. The story is known in the 12th chapter of Matthew, of the devil, proposing to return, and therefore taking unto him seven worse spirits. But why do you hear him, why do you obey him, whom you see shut out of heaven? Nevertheless, hereby is also signified that the devil was so fully vanquished by Christ, that he was also driven to forsake the place of the battle.

For the third member, as it were expounding the second, adds: and he was cast to the Earth. For those who are thrown to the ground are judged to be overcome. Therefore, a full victory and perfect conquest is signified. However, he was once most valiantly cast to the Earth. Of our Lord Jesus Christ, in the mystery of our redemption, and in the virtue of the same, is daily cast to the Earth, by the faithful. And just as the Devil has no permanent place in heaven nor in the elect, so verily he inhabits all earthly things, that is to say, men savoring the earth and contemning heavenly things. Yea, and we hear that his angels are cast out with him. For the Lord in the Gospel of Saint John, chapter 16, says: “In the world you have affliction, but be of good cheer, I have overcome the world.” And Saint John in his canonical epistle says, “You are of God, little children, and you have overcome them, for he is greater that is in you, than he is in the world.” And this is the victory that overcame the world, even your faith.

And by the way he explains what we should understand by the dragon, of whom he has spoken until now, namely, the old enemy of humankind. He depicts him with his titles, attributing four names to him, so that we may also better understand his nature and beware of that wicked murderer. First, he calls him the old Serpent. For at the beginning, by the serpent, he infected our first parents with the poison of death and sin, and by that, the entire world: as is seen in Genesis 3 and Romans 5. Therefore, I said at the beginning of this chapter that he is called a dragon. Afterward, he calls him the Devil, that is to say, a slanderer or a false accuser. For it immediately follows, which may explain this word: “the accuser of our brothers has been cast out,” etc. A fitting example of this thing is declared in the first and second chapters of Job. [illegible] signifies to accuse or blame, and [illegible] is an accusation, and [illegible] a crime or complaint.

Thirdly, he calls him Satanas, in the Hebrew word, namely an adversary, for he is in all things against God, and opposes himself and resists people in holy matters, if perhaps he might hinder or corrupt them. Lastly, he is called a deceiver, a betrayer, or one who subverts and betrays the entire world. For this the Lord attributes to him in John 8, for that he has been a liar from the beginning, and is the father, that is the fountain and origin of all lying, deception, errors, and seduction, and of all evil. For all errors and heresies, all deceptions, and all falsehoods, finally all kinds of evils, have flowed out of this most filthy wellspring. And who is he that hears these things, which will not abhor that vile beast? They must needs be stark mad who seek by all means to be in favor with that wicked spirit.

He should now consequently add the remainder of this account, namely how the Dragon persecutes and attacks the woman, and she again resists by fleeing and overcomes through Christ. But he suspends this narration for a little while, and places now a song of victory and triumph of the saints in heaven, of the angels and blessed souls. The sum of it is that Christ has overcome, and that the faithful overcome in Christ; and therefore even the heavens themselves, and all who dwell therein, must rejoice and sing. And I repeat that these things are interlaid in the dangerous Antichristian and Roman view, for a consolation, lest the saints should in those great dangers by reason of their natural infirmity be discouraged; but calling upon the name of Christ, should fight manfully, when they understand under whose banner they fight, and with whom they fight: truly, with one overcome under Christ’s standard. And when we hear that the Dragon’s force is broken, we shall think that the furies of either beast, as well the ten-horned as the two-horned, are weakened in the faith of Christ. This gives also no small courage in this conflict, that we see that the Dragon has no power over those who are sprinkled and purified with the blood of Christ, but over earthly and worldly men. And this triumph is heavenly. For voices are heard out of heaven, singing a merry note, to the intent that the rejoicing of the blessed spirits might have more authority, grace, and efficacy among the poor afflicted.

\index{Resurrection}They all with one voice maintain simply that salvation and power are now made perfect, for by the Lord’s death and resurrection, God has wrought power and made perfect the salvation promised to the fathers, namely while he trod down the serpent’s head, abolished sin and death, and restored life. Thus is the kingdom of God in this world established in the elect, while even by the power of Christ, the Prince of this world is cast out and overcome. For this is the reason we must rejoice, and to see the virtue and power of Christ revealed, and how salvation is made perfect: because, he says, through Christ, the Devil is cast down, that is to say, overcome and vanquished, so that he can no more accuse humankind before the judgment seat of God. To this belongs what Saint Paul wrote: Who shall accuse the elect of God? It is God that justifies; who is he that condemns? It is Christ, who died, yes, who rose again, who is also on the right hand of God, making intercession for us.

Moreover, the heavenly beings do not only preach the victory of Christ, but of all the faithful, which they obtain against Satan in the faith of Jesus Christ; that it may hereof at least appear what we should understand before by Michael and by his angels. And he emphasizes that Christians overcome Satan not by their own merits, force, or strength, but by the merit and grace of Christ. And they say, that is, the angels of Michael, overcame the Dragon by the blood of the Lamb. For inasmuch as the faithful are purified by the blood of Christ, Satan has nothing against them; but if they have the spirit and faith of Christ, they overcome the Devil also. So in times past, the destroyer had no power over those houses which were marked with the blood of the Lamb, Exodus 12. And he adds another thing for which the faithful overcame: for the word of the testimony of Christ, which is the gospel. Which because it is invincible and eternal, they overcome all things of this world, whoever abides in the lively and eternal word of truth. And even in the most true gospel, the Lord himself has promised that he will not forsake his own, and will fight for them. Therefore, the faithful must needs overcome. To these things is added the effect of Christ’s purifying: they loved not their life more than Christ; and therefore they have given it for Christ unto death, and so have overcome. For many are vanquished by this one thing, that they will not hazard their life for Christ.

For these great benefits of God, they now exhort the heavens themselves, and all the inhabitants of heaven, that is to say, they exhort one another to sing a joyful song. And that which the blessed saints say they do here, they teach the saints on earth to do also, and instruct them about what manner and sort they ought to be, who shall overcome Satan in battle—namely, purified by the blood of Christ, cleaving to the testimony of Jesus Christ, and contemptuous of their own lives, to whom it seems not grievous to die for Christ’s sake.

Finally, concerning the song’s ending, they declare that the Devil will reign and hold sway in earthly and fleshly men—those who truly deride godly things and value only worldly possessions, which are destined to perish. For the acquisition and preservation of these things, they will be willing to do anything, no matter how difficult. For Christ’s cause, they will endure nothing, neither action nor suffering. To these, they pronounce a terrible woe, namely, the curse of this present age and of the life to come.

But in him whom the Devil possesses in his kingdom, he also utters his malice against the elect, and that his great malice. For he rages most cruelly against the godly, and against godliness. He rages also most extremely against those his worshippers, whom he pollutes with all kinds of filthiness, and with all shame and reproach defiles.

\index{Judgment, Last}Again I suppose that same pertains to the comfort of the godly, that is spoken of a short time. For Satan, indeed, through Antichrist shall most cruelly rage against the church, but those days shall be shortened for the elect’s sake. By the way, it is noted also the wicked nature of Satan, which, knowing that the last judgment is at hand, wherein he must be thrown headlong into hell, thinks to requite and recompense the shortness of time with the cruelty of his wrath and devilish fury.

And thus far concerning the victory of Christ and his saints. Now follow, with less terror, yet horrible things concerning the war which the dragon most greedily and fiercely instigates against the mother of God. The Lord Jesus brings him under our feet. Amen. Amen.
